% ============================================================
% f2014_beijing_arts_q08_popcorn.tex
% 2014 北京卷（文科）第 8 题：可食用率 p 与加工时间 t 的三次实验数据
% 来源 PDF：普通高考/2014/2014北京文.pdf
% 原页：第 1 页（A3），印面 863，第 8 题题干上方偏右
% 三个数据点：(3, 0.7)、(4, 0.8)、(5, 0.5)
% 横轴 t 每 1 分钟 1.5cm；纵轴 p 每 0.1 为 1cm。
% ============================================================
\begin{tikzpicture}[>=Stealth, line width=0.6pt, font=\small]
  % 坐标轴
  \draw[->] (0,0) -- (0,8.55) node[above] {$p$};
  \draw[->] (0,0) -- (8.8,0) node[below left=-1pt] {$t$};
  \node[below left] at (0,0) {$O$};

  % 纵轴刻度值
  \node[left] at (0,5) {$0.5$};
  \node[left] at (0,7) {$0.7$};
  \node[left] at (0,8) {$0.8$};

  % 横轴刻度值
  \node[below] at (4.5,0) {$3$};
  \node[below] at (6.0,0) {$4$};
  \node[below] at (7.5,0) {$5$};

  % 参考虚线
  \draw[dashed, line width=0.45pt] (0,7.0) -- (4.5,7.0);
  \draw[dashed, line width=0.45pt] (0,8.0) -- (6.0,8.0);
  \draw[dashed, line width=0.45pt] (0,5.0) -- (7.5,5.0);
  \draw[dashed, line width=0.45pt] (4.5,0) -- (4.5,7.0);
  \draw[dashed, line width=0.45pt] (6.0,0) -- (6.0,8.0);
  \draw[dashed, line width=0.45pt] (7.5,0) -- (7.5,5.0);

  % 三个数据点
  \fill (4.5,7.0) circle (1.5pt);
  \fill (6.0,8.0) circle (1.5pt);
  \fill (7.5,5.0) circle (1.5pt);
\end{tikzpicture}
